% Worked Examples: Concept 2.5.5 - Zero Frequency Gain, Poles, and Zeros
\documentclass[11pt,a4paper]{article}
\usepackage{worked-examples-template}

\title{\textbf{Worked Examples}\\[0.5em]
\large Concept 2.5.5: Zero Frequency Gain, Poles, and Zeros\\[0.3em]
\normalsize \coursesubtitle{} -- \coursetitle}
\author{Assoc.\ Prof.\ William Robertson \and Dr Sean McGowan}
\date{\today}

\begin{document}

\maketitle
\tableofcontents
\newpage

\section{Concept 2.5.5: DC Gain}

\subsection{Example 1: Computing DC Gain and Steady-State Response}

\paragraph{Problem:} For the transfer function:
\begin{align}
G(s) = \frac{10(s+2)}{(s+1)(s+5)}
\end{align}

Find: (a) the DC gain, and (b) the steady-state output for a unit step input.

\paragraph{Solution:}

\textbf{Part (a): DC Gain}

The DC gain is the value of the transfer function at $s = 0$:
\begin{align}
G(0) = \frac{10(0+2)}{(0+1)(0+5)} = \frac{20}{5} = 4
\end{align}

\textbf{Part (b): Steady-State Response to Unit Step}

For a unit step input, $U(s) = \frac{1}{s}$, the output is:
\begin{align}
Y(s) = G(s)U(s) = \frac{10(s+2)}{(s+1)(s+5)} \cdot \frac{1}{s}
\end{align}

Using the final value theorem:
\begin{align}
y_{ss} = \lim_{t \to \infty} y(t) = \lim_{s \to 0} sY(s)
\end{align}

Compute:
\begin{align}
y_{ss} &= \lim_{s \to 0} s \cdot \frac{10(s+2)}{s(s+1)(s+5)} \\
&= \lim_{s \to 0} \frac{10(s+2)}{(s+1)(s+5)} \\
&= \frac{10(2)}{(1)(5)} = \frac{20}{5} = 4
\end{align}

\textbf{Result:} 
\begin{itemize}
\item DC gain: $G(0) = 4$
\item Steady-state step response: $y_{ss} = 4$
\end{itemize}

\textbf{Interpretation:} The DC gain equals the steady-state response to a unit step input, which is a fundamental property of stable systems. An input step of magnitude 1 results in an output that settles to a value of 4.
\end{document}
